% !TEX root = ../../tree_notes.tex
% =====================================================================
%  16.5  Distances, Paths and Depth Buckets             sec:c11-paths
%  (Renamed from 04_paths_lca.tex for the 5-section layout)
%  Exports: thm:c11-distf  prop:c11-onpath  prop:c11-interval
%           ex:c11-dists  ex:c11-cousinsrun  lst:c11-cousins
%           prob:c11-cousins  lst:c11-208e-soln
% =====================================================================
\section{Distances, Paths and Depth Buckets}\label{sec:c11-paths}

\begin{previouslybox}[Four questions from 16.1--16.4]
\begin{enumerate}[nosep, leftmargin=1.8em]
  \item From \(\exref{c11-census}\): \(\lca(8,10) = 1\).  \emph{(How many
        edges from~8 to~10 via the hinge?)}
  \item From \(\secref{c11-euler}\): the distance formula
        \(\dist(u,v) = \depth(u)+\depth(v)-2\depth(\lca)\) is one
        arithmetic expression given the LCA.
  \item From \(\secref{c11-lifting}\): the \(k\)-th ancestor gives route
        positions (section 16.4's ``\(k\)-th node on a route'').
  \item From Ch.~13 (\(\secref{c10-distance}\)): the identity
        \(\dist(s,u)+\dist(u,v)=\dist(s,v)\) iff \(u\) lies on the route
        \(s\to v\).  \emph{(That is today's \propref{c11-onpath}.)}
\end{enumerate}
\end{previouslybox}

\begin{topicintro}
The LCA is a \emph{measurement hub}: once you have \(\lca(u,v)\), the
distance follows for free (\thmref{c11-distf}); whether a third node
lies on the route follows from a triangle equality
(\propref{c11-onpath}); whether one node is an ancestor of another
follows from the \(\tin/\tout\) interval test (\propref{c11-interval},
re-stated from Ch.~5 for convenience).  Finally, \emph{per-depth
buckets} -- for each depth~\(d\), a sorted list of \(\tin\) values of all
nodes at depth~\(d\) -- answer ``how many descendants of the \(p\)-th
ancestor of~\(v\) sit at the same depth as~\(v\)?'' in \(O(\log n)\),
the key to CF~208E (Blood Cousins).  This section closes the static-tool
chapter: everything here is queryable with \(O(n\log n)\) preprocessing
and \(O(\log n)\) per question, and the next chapter (Ch.~17 DSU-on-tree)
builds on these tools.

\textbf{After this section you can}
\begin{itemize}[nosep]
  \item prove the distance-through-hinge theorem and compute any
        pairwise distance by hand on~\(T\); \hfill\lv{3}
  \item test whether~\(w\) is on the route \(u\)-\(v\) using triangle
        equality (\propref{c11-onpath}); \hfill\lv{3}
  \item apply the \(\tin/\tout\) subtree interval test
        (\propref{c11-interval}); \hfill\lv{2}
  \item build per-depth \(\tin\)-bucket arrays and binary-search a
        subtree count in \(O(\log n)\); \hfill\lv{4}
  \item solve CF~208E (Blood Cousins) end-to-end. \hfill\lv{4}--\lv{5}
\end{itemize}
\end{topicintro}

% ---------------------------------------------------------------
\subsection{Distance through the Hinge}
% ---------------------------------------------------------------

\begin{theorem}{Distance formula}{c11-distf}
For any nodes \(u, v\) in a tree rooted at~\(r\) (unweighted edges):
\[
  \dist(u,v) \;=\; \depth(u) + \depth(v) - 2\,\depth(\lca(u,v)).
\]
For edge-weighted trees, replace \(\depth\) by the weighted
\(\dist(r,\cdot)\) and the formula holds identically.
\end{theorem}

\begin{proofbox}
The unique route \(u\to v\) passes through \(\ell = \lca(u,v)\).
The segment \(u\to\ell\) has length \(\depth(u)-\depth(\ell)\) (one edge
per ancestor step).  Symmetrically \(\depth(v)-\depth(\ell)\).  The two
segments share only \(\ell\), so the total is their sum.
\end{proofbox}

\begin{example}{Distance table on T}{c11-dists}
Using \(\depth = (0,1,1,2,2,2,2,3,3,4)\) and the census for LCAs:
\[
\renewcommand{\arraystretch}{1.2}
\begin{array}{c|c|c|c|c}
u & v & \lca(u,v) & \depth(u)+\depth(v)-2\depth(\lca) & \dist(u,v) \\ \hline
8  & 10 & 1 & 3+4-2(0) & 7\\
4  & 8  & 2 & 2+3-2(1) & 3\\
9  & 10 & 9 & 3+4-2(3) & 1\\
2  & 10 & 1 & 1+4-2(0) & 5\\
5  & 6  & 1 & 2+2-2(0) & 4\\
4  & 5  & 2 & 2+2-2(1) & 2\\
7  & 10 & 7 & 2+4-2(2) & 2\\
6  & 10 & 3 & 2+4-2(1) & 4
\end{array}
\]
Every row is verifiable by counting edges on the drawn tree.
\end{example}

% ---------------------------------------------------------------
\subsection{On-Path Test (Triangle Equality)}
% ---------------------------------------------------------------

\begin{proposition}{\(w\) lies on the route \(u\)-\(v\)}{c11-onpath}
For three nodes \(u, v, w\) in a tree:
\[
  w \in \text{path}(u,v) \iff \dist(u,w)+\dist(w,v) = \dist(u,v).
\]
\end{proposition}

\begin{proofbox}
The forward direction: if \(w\) is on the route, removing \(w\) splits it
into two segments of lengths \(\dist(u,w)\) and \(\dist(w,v)\) that
exactly tile \(\dist(u,v)\).  The reverse: if the equality holds, the
unique route \(u\to v\) must pass through \(w\) because the shortest path
metric in a tree satisfies the triangle inequality with equality iff the
midpoint is on the geodesic.
\end{proofbox}

\begin{tracebox}[On-path test on T: is 7 on path 8 to 10?]
\(\dist(8,7) = 3+2-2(0) = 5\); \(\dist(7,10) = 2+4-2(2) = 2\);
\(\dist(8,10) = 7\).  Check: \(5+2 = 7 = \dist(8,10)\).
\(\Rightarrow\) 7 IS on the path \(8\to10\).  \checkmark

Is~6 on path \(8\to10\)?  \(\dist(8,6)=3+2-0=5\);
\(\dist(6,10)=2+4-2(1)=4\); \(5+4=9\ne 7\).  NO.  \checkmark
\end{tracebox}

% ---------------------------------------------------------------
\subsection{Ancestry via tin/tout (Back-cite Ch.~5)}
% ---------------------------------------------------------------

\begin{proposition}{Subtree-as-interval test}{c11-interval}
Under the entry-only DFS of Chapter~5 (\(\chref{c24}\)), with
\(\tout(v)\) = the last \(\tin\) stamp inside \(v\)'s subtree:
\[
  u \preceq v \iff \tin(u) \le \tin(v) \le \tout(u).
\]
Cost \(O(1)\) per test after \(O(n)\) preprocessing.
\end{proposition}

\noindent\textbf{Note.}  This is the \emph{same} \(\tin/\tout\) array
printed in \figref{c11-tree}, not the \(2n\!-\!1\) tour of
\(\secref{c11-euler}\).  Both coexist; they serve different purposes.

\begin{tracebox}[tin/tout test on T]
\(\tin = (1,2,6,3,4,7,8,5,9,10)\) for nodes 1..10.
\(\tout = (10,5,10,3,5,7,10,5,10,10)\).

Is \(7 \preceq 10\)?  \(\tin(7)=8 \le \tin(10)=10 \le \tout(7)=10\).  YES.  \checkmark
Is \(9 \preceq 6\)?  \(\tin(9)=9 \le \tin(6)=7\)?  \(9 > 7\).  NO.  \checkmark
Is \(2 \preceq 8\)?  \(\tin(2)=2 \le \tin(8)=5 \le \tout(2)=5\).  YES.  \checkmark
\end{tracebox}

% ---------------------------------------------------------------
\subsection{Per-Depth Buckets}
% ---------------------------------------------------------------

For each depth~\(d\), maintain a sorted array \(\mathit{bucket}[d]\)
containing \(\tin(v)\) for every \(v\) with \(\depth(v) = d\).  To count
nodes at depth~\(d\) that are in the subtree of~\(w\) (i.e.,
\(\tin(w) \le \tin(v) \le \tout(w)\) and \(\depth(v)=d\)):
\[
  \text{count} = \texttt{upper\_bound}(\mathit{bucket}[d],\;\tout[w])
  \;-\; \texttt{lower\_bound}(\mathit{bucket}[d],\;\tin[w]).
\]
Cost \(O(\log n)\) per query (binary search in a sorted array).

\begin{tracebox}[Depth buckets on T (sorted by tin)]
\[
\renewcommand{\arraystretch}{1.2}
\begin{array}{c|l|l}
d & \text{nodes at depth } d & \text{bucket[d] (tin values sorted)} \\ \hline
0 & \{1\} & [1] \\
1 & \{2,\;3\} & [2, 6] \\
2 & \{4,\;5,\;6,\;7\} & [3, 4, 7, 8] \\
3 & \{8,\;9\} & [5, 9] \\
4 & \{10\} & [10]
\end{array}
\]
Example: count nodes at depth~2 in subtree of~2
(\(\tin(2)=2\), \(\tout(2)=5\)).  Search \(\mathit{bucket}[2]=[3,4,7,8]\)
for values in \([2,5]\): \(\{3,4\}\), count~=~2 (nodes 4 and~5).  \checkmark
\end{tracebox}

% ---------------------------------------------------------------
\subsection{Problem Lab}
% ---------------------------------------------------------------

\solvedproblem{c11-cousins}{Blood Cousins}{CF 208E}
  {https://codeforces.com/problemset/problem/208/E}{CF 2100 L5}
  {binary lifting, k-th ancestor, depth buckets, subtree tin/tout}

\begin{problemstatement}[CF 208E Blood Cousins]
A \emph{forest} of boys is given by a father array of~$n$ integers
($0$~marks a root --- the input need not be one tree); then $q$~queries
$(v, p)$ with $p > 0$: how many boys $u \ne v$ have the same
$p$-th ancestor as~$v$?  (Equivalently: let $w$ be the $p$-th
ancestor of~$v$; count nodes at depth $\depth(v)$ inside subtree($w$),
minus~$v$ itself.)  If $p > \depth(v)$ the ancestor doesn't exist;
the answer is~$0$.

\textbf{Input.} $n$, then a line of $n$ integers $a_1, \ldots, a_n$
($a_i$ is the parent of node~$i$; $a_i = 0$ means node~$i$ is a root).
Then $q$, then $q$ lines each with $v$ and~$p$.
$1 \le n, q \le 10^5$.

\textbf{Example.}
\begin{verbatim}
Input:
6
0 1 1 0 4 4
7
1 1
1 2
2 1
2 2
4 1
5 1
6 1

Output:
0 0 1 0 0 1 1
\end{verbatim}
\end{problemstatement}

\begin{attemptgate}
\peek{Identify the root (the node with parent 0), build depth and
\(\tin/\tout\), then for each query climb \(p\) steps via lifting.}
\nudge{After finding \(w = \text{\(p\)-th ancestor of \(v\)}\), count
nodes at \(\depth(v)\) in \(w\)'s subtree using the bucket + binary search.}
\fullhint{Answer = (upper\_bound - lower\_bound in bucket[depth(v)]
restricted to [tin(w), tout(w)]) - 1.}
\end{attemptgate}

\solstep{Understanding}
The query is ``how many people are at my generation (same depth) and
share the \(p\)-th ancestor (same great\(\ldots\)grandparent)?''  The
translation: count same-depth nodes in the ancestor's subtree, minus
yourself.  Three operations: \(k\)-th ancestor (lifting), depth lookup,
range count in a bucket (binary search).

\solstep{Key Observations}
\begin{observation}
The input is a \emph{forest} --- multiple nodes may have parent~0 (the
sample has roots~1 and~4).  A super-root~0 glues all real roots as
children; depths shift by~1 but the algorithm is unchanged.
\end{observation}
\begin{observation}
If $p \ge \depth(v)$ (depth of~$v$ from its real root) the climb
overshoots; answer~0.  Guard before calling \texttt{kth\_ancestor}.
\end{observation}

\solstep{Think It Through}
\begin{thinkbox}
\deadend{DFS from \(v\) counting same-depth nodes}
{\(O(n)\) per query; \(O(nq)\) total.  TLE.}
\deadend{Precompute a 2D array of subtree depth histograms}
{Space \(O(n^2)\).  MLE.}
\aha{Bucket nodes by depth (sorted by \(\tin\)).  Per query: find \(w\)
(\(p\)-th ancestor, \(O(\log n)\)), then binary-search bucket[depth(v)]
for the interval [tin(w), tout(w)].  \(O(\log n)\) per query, \(O(n)\) space.}
\end{thinkbox}

\solstep{Implementation}
\begin{lstlisting}[caption={CF 208E full solution (super-root 0; lifting + depth buckets)}, label={lst:c11-208e-soln}]
#include <bits/stdc++.h>
using namespace std;
const int LOG=17, MX=100002;
vector<int> adj[MX], bucket[MX];
int up[MX][LOG], dep[MX], tin[MX], tout[MX], tmr;

void dfs(int v, int p) {
    tin[v]=++tmr; up[v][0]=p;
    for(int j=1;j<LOG;j++) up[v][j]=up[up[v][j-1]][j-1];
    bucket[dep[v]].push_back(tin[v]);   // sorted by construction
    for(int c:adj[v]) if(c!=p){dep[c]=dep[v]+1; dfs(c,v);}
    tout[v]=tmr;
}
int kth(int v,int k){
    for(int j=0;j<LOG;j++) if(k&(1<<j)) v=up[v][j];
    return v;
}
int main(){
    int n; scanf("%d",&n);
    dep[0]=-1;                           // super-root 0
    for(int i=1;i<=n;i++){
        int p; scanf("%d",&p);
        if(p==0) adj[0].push_back(i);    // root -> child of super-root
        else adj[p].push_back(i);
    }
    up[0][0]=0; for(int j=1;j<LOG;j++) up[0][j]=0;  // sentinel row
    dep[0]=0; dfs(0,0);                   // one pass over the forest
    // dep of real roots = 1 (one below super-root); adjust if needed
    int q; scanf("%d",&q);
    while(q--){
        int v,p; scanf("%d%d",&v,&p);
        if(p>dep[v]-1){printf("0 ");continue;} // overshoot guard
        int w=kth(v,p);
        int d=dep[v];
        int lo=lower_bound(bucket[d].begin(),bucket[d].end(),tin[w])-bucket[d].begin();
        int hi=upper_bound(bucket[d].begin(),bucket[d].end(),tout[w])-bucket[d].begin();
        printf("%d ",hi-lo-1);
    }
}
\end{lstlisting}

\solstep{Trace on the Official Sample}
\begin{tracebox}[CF 208E: forest 0~1~1~0~4~4; all 7 queries verified]
Forest structure: root~1 has children~{2,3}; root~4 has children~{5,6}.
Depths (from each real root): $\depth = (0,1,1,0,1,1)$ for nodes 1..6.

Q1\;$(v\!=\!1,\, p\!=\!1)$: $p\!=\!1 > \depth(1)\!=\!0$; overshoot
$\Rightarrow$~\textbf{0}. \checkmark

Q2\;$(1,\,2)$: $p\!=\!2 > \depth(1)\!=\!0$; overshoot
$\Rightarrow$~\textbf{0}. \checkmark

Q3\;$(2,\,1)$: 1st~anc of~2 is~1.  Depth~1 nodes in subtree(1)
$= \{2,3\}$; count~$2 - 1 = $~\textbf{1}. \checkmark

Q4\;$(2,\,2)$: $p\!=\!2 > \depth(2)\!=\!1$; overshoot
$\Rightarrow$~\textbf{0}. \checkmark

Q5\;$(4,\,1)$: $p\!=\!1 > \depth(4)\!=\!0$; overshoot
$\Rightarrow$~\textbf{0}. \checkmark

Q6\;$(5,\,1)$: 1st~anc of~5 is~4.  Depth~1 nodes in subtree(4)
$= \{5,6\}$; count~$2 - 1 = $~\textbf{1}. \checkmark

Q7\;$(6,\,1)$: 1st~anc of~6 is~4.  Depth~1 nodes in subtree(4)
$= \{5,6\}$; count~$2 - 1 = $~\textbf{1}. \checkmark

Output: \texttt{0 0 1 0 0 1 1}. \checkmark
\end{tracebox}

\begin{example}{Cousin-count queries on T}{c11-cousinsrun}
On \(T\) (depths \(d = (0,1,1,2,2,2,2,3,3,4)\)):
\[
\renewcommand{\arraystretch}{1.2}
\begin{array}{c|c|c|l|c}
v & p & w = \text{\(p\)-th anc of \(v\)}
  & \text{depth}=\depth(v)\text{ nodes in subtree}(w) & \text{ans} \\ \hline
8 & 1 & 5  & \text{depth~3 in sub(5)=\{5,8\}: \{8\}} & 0 \\
9 & 1 & 7  & \text{depth~3 in sub(7)=\{7,9,10\}: \{9\}} & 0 \\
4 & 1 & 2  & \text{depth~2 in sub(2)=\{2,4,5,8\}: \{4,5\}} & 1 \\
4 & 2 & 1  & \text{depth~2 in sub(1)=all: \{4,5,6,7\}} & 3 \\
8 & 2 & 2  & \text{depth~3 in sub(2)=\{2,4,5,8\}: \{8\}} & 0 \\
10& 2 & 7  & \text{depth~4 in sub(7)=\{7,9,10\}: \{10\}} & 0 \\
10& 3 & 3  & \text{depth~4 in sub(3)=\{3,6,7,9,10\}: \{10\}} & 0 \\
6 & 2 & 1  & \text{depth~2 in sub(1)=all: \{4,5,6,7\}} & 3 \\
1 & 1 & \text{overshoot} & - & 0 \\
\end{array}
\]
All hand-verified.  The last row is the overshoot guard (\(p=1 >
\depth(1)=0\): answer~0).
\end{example}

\solstep{Complexity}
\begin{complexitybox}
Preprocessing: \(O(n\log n)\) (lifting table) + \(O(n)\) (bucket build,
sorted since \(\tin\) increases during DFS in depth order).
Per query: \(O(\log n)\) (climb) + \(O(\log n)\) (binary search).
Total \(O((n+q)\log n)\).  \textbf{AC.}
\end{complexitybox}

% ---------------------------------------------------------------
\subsection{Common Pitfalls}
% ---------------------------------------------------------------

\begin{warningbox}[Depth-0 nodes and the overshoot guard]
If \(p > \depth(v)\), the \(p\)-th ancestor is the sentinel~0 (or
undefined).  Querying \(\mathit{bucket}[\depth(v)]\) for the interval
\([\tin(0), \tout(0)]\) would count everything (or crash).  Always
check \(p > \depth(v)\) FIRST and return~0.
\end{warningbox}

\begin{warningbox}[The -1 (exclude v itself) must come after the count]
If \(v\) is in \(w\)'s subtree (always true when \(w\) is an ancestor
of~\(v\)), the binary-search count \(\ge 1\).  Forgetting the \(-1\)
counts \(v\) as its own ``cousin''.  The answer is the range size minus~1.
\end{warningbox}

\begin{warningbox}[Mixing tin/tout conventions]
The bucket stores \(\tin(v)\) values.  The interval test for subtree is
\(\tin(w) \le x \le \tout(w)\) (last-stamp convention from Ch.~5).  If
your DFS uses the ``exit tick'' convention (\(\tout(w) =\) timestamp at
the moment of exiting), the interval upper-bound is off by~1 and you
over-count.  Pin ONE convention.
\end{warningbox}

% ---------------------------------------------------------------
\subsection{Variations and Twists}
% ---------------------------------------------------------------

\begin{itemize}[nosep]
  \item \textbf{LC 236 (Ch.~1 revisit).}  The binary-tree LCA from
        LeetCode 236 is \(\exref{c11-census}\) restricted to a
        \emph{binary} tree where the ``walk'' fits in a post-order
        recursion.  Prose revisit only (the URL lives in Ch.~1's mixed
        set).  \hfill\lv{2}
  \item \textbf{CSES 1136 Tree Distances II.}  A ``sum of distances
        from~\(v\) to all other nodes'' problem that combines Ch.~15's
        rerooting with \thmref{c11-distf}: the per-root sum is one
        reroot-every-node quantity.  Prose revisit (URL owned by Ch.~12
        notes).  \hfill\lv{5}
  \item \textbf{On-path count via tin/tout.}  Given the route \(u\to v\),
        the set \(\{w : \text{on-path}\}\) is not a subtree interval;
        but \(\{w : \text{on-path and in subtree}(w_0)\}\) is.  Useful in
        some CF~960-family queries.  \hfill\lv{5}
\end{itemize}

% ---------------------------------------------------------------
\subsection{Practice Corner}
% ---------------------------------------------------------------

\begin{practiceset}
  \practice{Find Distance in a Binary Tree}{LC 1740}
    {https://leetcode.com/problems/find-distance-in-a-binary-tree/}
    {LC Medium L3}
    {distance between two valued nodes: locate both, take their hinge,
     apply \(\dist(u,v)=\depth u+\depth v-2\depth\ell\).  One DFS records
     parent/depth and LC~236's LCA gives the hinge -- the binary-tree
     dress-up of \thmref{c11-distf}.}
\end{practiceset}

\paragraph{Back-cites.}
Ch.~5 (\(\chref{c24}\)) -- owns \(\tin/\tout\) flattening; we use
\propref{c11-interval} as a restatement, never re-derive.
Ch.~13 (\(\secref{c10-distance}\)) -- the triangle-equality identity used
by \propref{c11-onpath}; cited, not re-proved.
Ch.~15 (\(\chref{c20}\)) -- rerooting converts per-root quantities; CSES
1136 needs the two-DFS template + distance.

% ---------------------------------------------------------------
\begin{recallbox}[Section 16.5 in 10 seconds]
\thmref{c11-distf}: \(\dist(u,v) = \depth(u)+\depth(v)-2\depth(\lca)\).
\propref{c11-onpath}: \(w\) on route \(u\)-\(v\) iff
\(\dist(u,w)+\dist(w,v)=\dist(u,v)\).
\propref{c11-interval}: \(u\preceq v\) iff \(\tin(u)\le\tin(v)\le\tout(u)\)
(Ch.~5).  Per-depth buckets: \(\mathit{bucket}[d]\) sorted tins; subtree
count via binary search.  CF~208E: \(k\)-th ancestor + bucket range --
count depth-\(d\) nodes in \(w\)'s subtree, minus~1.
\selftest{
  \item \(\dist(8,10)\) on~\(T\)?  (7.)  Show the arithmetic.
  \item State \propref{c11-onpath} and test node~6 on path \(8\to10\).
  \item Bucket[2] on~\(T\)?  (\([3,4,7,8]\).)
  \item CF~208E query (4,2): what is \(w\)?  What count minus 1?  (1, 4-1=3.)
  \item What is the overshoot guard condition?
}
\end{recallbox}
